\begin{comment}
\begin{center}
\Large\textbf{PREFACE}
\end{center}
\vspace{2ex}
%Tulis kata pengantar di sini

Alhamdulillah, All praise and gratitude are dedicated to the One True Almighty \textbf{Allah Subhanahu Wa Ta'ala (SWT)} for His blessings, mercy, and kindness, which have enabled the author to complete the Doctoral Dissertation entitled \textit{"Cardiac Segmentation and Scar Quantification on CMR Images: A Deep Learning Approach Toward Myocardial Infarction (MI) Assessment"} This dissertation is prepared as part of the requirements to obtain a Doctoral degree in the Department of Electrical Engineering, Faculty of Intelligent Electrical and Informatics Technology, Institut Teknologi Sepuluh Nopember, Surabaya.

The author fully realizes that the completion of this dissertation would not have been possible without the support, guidance, and encouragement of many individuals and institutions. On this occasion, the author wishes to convey her deepest gratitude and sincere appreciation to:

\begin{enumerate}
	\item Prof. Dr. Ir. Mauridhi Hery Purnomo, M.Eng., the principal supervisor, for his prayers, patience, continuous motivation, insightful guidance, and invaluable advice throughout the doctoral journey. The author is truly grateful for the privilege of being mentored by him, which has provided not only academic enrichment but also meaningful intellectual and practical experiences, both nationally and internationally.
	
	\item Prof. Dr. I Ketut Eddy Purnama, S.T., M.T., the first co-supervisor. His academic rigor, non-academic encouragement, technical guidance, and unique perspectives brought clarity and depth to this work, particularly in the field of Medical Imaging. Above all, his mentorship and thoughtful support were a constant source of motivation. %for his invaluable academic guidance and non-academic encouragement, technical support, and enlightening perspectives, which have significantly contributed to the clarity, depth, and overall quality of this dissertation particularly in the field Medical Image. The author extends profound gratitude for his mentorship and thoughtful encouragement throughout the research process.
	
	\item Dr. Eko Mulyanto Yuniarno, S.T., M.T., the second co-supervisor. His academic advice, moral encouragement, technical support, warm hospitality, and imaginative thinking greatly enriched this dissertation, especially in Computer Vision. His mentorship and genuine care throughout the doctoral studies will always be remembered. %for his academic and non-academic encouragement,technical support, generous hospitality, and imaginative perspectives which have significantly contributed to the clarity, depth, and overall quality of this dissertation particularly in the field Computer Vision. The author extends profound gratitude for his mentorship and thoughtful encouragement throughout the doctoral studies.
	
	\item Prof. Dr. Eng. Masayoshi Aritsugi, the International co-supervisor and examiner, for his invaluable mentorship and for graciously providing the opportunity to conduct research in his laboratory at Kumamoto University through the exchange program. His expert technical guidance, rigorous feedback, and continuous support throughout the scientific writing process, from initial drafting to successful publication, brought exceptional clarity and depth to this work. Above all, his thoughtful encouragement and unique perspectives served as a constant source of inspiration and motivation.
	
	\item The Center for Higher Education Funding and Assessment (PPAPT) of the Ministry of Higher Education, Science, and Technology, Indonesia, and the Indonesian Endowment Fund for Education (LPDP), for their generous financial support through the Beasiswa Pendidikan Indonesia (BPI) program.
	
	\item All leaders, lecturers, and staff of the Department of Electrical Engineering, Faculty of Intelligent Electrical and Informatics Technology, Institut Teknologi Sepuluh Nopember, for their cooperation, guidance, and assistance throughout the author’s academic journey.
	
	\item All leaders, colleagues, and staff of Politeknik Negeri Jakarta, particularly from the Department of Electrical Engineering, for their continuous support and the valuable opportunities that have significantly contributed to the author’s professional growth.
	
	\item The author's late parents, Koendari Soelendar and Larasati. Deep remembrance and eternal gratitude are dedicated to her late \textit{Bapak} and late \textit{Ibu} for their unwavering love, prayers, enduring faith, and unconditional sacrifices, which remain the author's ultimate anchor and guiding light throughout this journey.
	
	\item The author's beloved big brother, Nayoko, and his wife, Yuliana Anggraini, together with their children, Aida Farah Angeline, Nawang Wulan Yulfamaretha, and Danang Sutawijaya; and her beloved younger sister, Mardiarini, and her husband, Jatmiko Nugroho, as well as their children, Fawnia Anjani Nugroho and Rayyan Aufa Nugroho, for their endless prayers, encouragement, and unwavering support, which served as a constant source of strength and comfort.

	%\item The author’s beloved late mother, Larasati, and beloved late father, Koendari Soelendar, for their %unwavering prayers, love, encouragement, and moral support, which have been of great significance throughout %this journey. Also, extends heartfelt thanks to my beloved big brother, Nayoko, and his wife, Yuliana %Anggraini, along with their children, Aida Farah Angeline, Nawang Wulan Yulfamaretha, and Danang Sutawijaya; %and to my beloved little sister, Mardiarini, and her husband, Jatmiko Nugroho, along with their children, %Fawnia Anjani Nugroho and Rayyan Aufa Nugroho, for their prayers, encouragement, and unconditional support, %which have been the greatest source of strength and motivation throughout this journey.

	\item All colleagues of the VISIKOM Laboratory, for cultivating a warm, supportive, and collegial environment during this academic endeavor. The author looks forward to continuing fruitful scientific collaborations in the future.
	
	\item All individuals whose names cannot be mentioned one by one, yet who have provided prayers, support, and inspiration in countless ways.
\end{enumerate}

The author extends her profound gratitude to all who have contributed, with the prayer that Allah SWT may bless them with health, happiness, and continued success. The author also humbly acknowledges that this dissertation may still contain limitations and warmly welcomes constructive feedback and suggestions for improvement. It is the author’s sincere hope that this work will provide meaningful insights and contribute positively to both academic discourse and practical applications. Aamiin yaa rabbal ’aalamiin.
	
	
	
	\vspace{26pt}
	\begin{flushright}
		\begin{tabular}[b]{c}
			Surabaya, ??? 2026
			\\
			\\
			\\
			Author
		\end{tabular}
	\end{flushright}
\end{comment}

\begin{center}
	\large\textbf{PREFACE}
\end{center}
\vspace{1ex}

Alhamdulillah, all praise and gratitude belong to \textbf{Allah Subhanahu Wa Ta'ala (SWT)} for His blessings and mercy, enabling the author to complete this Doctoral Dissertation entitled \textit{"Cardiac Segmentation and Scar Quantification on CMR Images: A Deep Learning Approach Toward Myocardial Infarction (MI) Assessment"}. This dissertation is submitted in partial fulfillment of the requirements for the Doctoral degree at the Department of Electrical Engineering, Faculty of Intelligent Electrical and Informatics Technology, Institut Teknologi Sepuluh Nopember, Surabaya.

The author expresses her deepest appreciation and gratitude to:

\begin{enumerate}
	\setlength{\itemsep}{1pt}
	\setlength{\parskip}{0pt}
	\setlength{\parsep}{0pt}
	
	\item Prof. Dr. Ir. Mauridhi Hery Purnomo, M.Eng., the principal supervisor, for his prayers, patience, continuous motivation, insightful guidance, and invaluable advice throughout the doctoral journey. The author is truly grateful for the privilege of being mentored by him, which has provided not only academic enrichment but also meaningful intellectual and practical experiences, both nationally and internationally.

	\item Prof. Dr. I Ketut Eddy Purnama, S.T., M.T., the first co-supervisor. His academic rigor, non-academic encouragement, technical guidance, and unique perspectives brought clarity and depth to this work, particularly in the field of Medical Imaging. Above all, his mentorship and thoughtful support were a constant source of motivation. 

	\item Dr. Eko Mulyanto Yuniarno, S.T., M.T., the second co-supervisor. His academic advice, moral encouragement, technical support, warm hospitality, and imaginative thinking greatly enriched this dissertation, especially in Computer Vision. His mentorship and genuine care throughout the doctoral studies will always be remembered. 

	\item Prof. Dr. Eng. Masayoshi Aritsugi, the International examiner, for his invaluable mentorship and for graciously providing the opportunity to conduct research in his laboratory at Kumamoto University through the exchange program. His expert technical guidance, rigorous feedback, and continuous support throughout the scientific writing process, from initial drafting to successful publication, brought exceptional clarity and depth to this work. Above all, his thoughtful encouragement and unique perspectives served as a constant source of inspiration and motivation.
	
	\item The Center for Higher Education Funding and Assessment (PPAPT) of the Ministry of Higher Education, Science, and Technology, Indonesia, and LPDP for financial support through the Beasiswa Pendidikan Indonesia (BPI) program.
	
	\item Leaders, lecturers, and staff of the Department of Electrical Engineering, ITS, for their continuous support and assistance.
	
	\item Leaders, colleagues, and staff of Politeknik Negeri Jakarta, especially the Department of Electrical Engineering and the Study Program of Instrumentation and Industrial Control , for their valuable support and professional opportunities.
	
	\item Her late parents, Koendari Soelendar and Larasati. Eternal gratitude is dedicated to her late \textit{Bapak} and late \textit{Ibu} for their unconditional love, prayers, and sacrifices, which remain her ultimate guiding light.
	
	\item Her big brother, Nayoko, and his wife, Yuliana Anggraini, together with their children, Aida Farah Angeline, Nawang Wulan Yulfamaretha, and Danang Sutawijaya; and her younger sister, Mardiarini, and her husband, Jatmiko Nugroho, as well as their children, Fawnia Anjani Nugroho and Rayyan Aufa Nugroho, for their endless prayers, encouragement, and unwavering support, which served as a constant source of strength and comfort.
	
	\item Colleagues at VISIKOM Laboratory for creating a warm, supportive, and inspiring research environment. Looking forward to continuing fruitful scientific collaborations in the future.
	
	\item All individuals whose names cannot be mentioned individually, but whose support and prayers have been invaluable.
\end{enumerate}

The author extends her profound gratitude to all who have contributed, with the prayer that Allah SWT may bless them with health, happiness, and continued success. The author also humbly acknowledges that this dissertation may still contain limitations and warmly welcomes constructive feedback and suggestions for improvement. It is the author’s sincere hope that this work will provide meaningful insights and contribute positively to both academic discourse and practical applications. Aamiin yaa rabbal ’aalamiin.

\nopagebreak
\vspace{2ex}
\begin{flushright}
	\begin{minipage}{5cm}
		\centering
		Surabaya, \today \\[4ex]
		Author
	\end{minipage}
\end{flushright}